\documentclass[11pt]{article}
\usepackage[margin=1.05in]{geometry}
\usepackage{amsmath,amssymb,amsthm,mathtools}
\usepackage{enumitem}
\usepackage[colorlinks,linkcolor=blue!55!black,citecolor=blue!55!black]{hyperref}

\makeatletter
%% protocol-header.tex -- PROTOCOL §2.3 page-1 block, as a TEMPLATE not an instance.
%%
%% PROTOCOL §2.3: "Every PDF must state, on its first page: author, date,
%% recipient, a title, and the commit hash of the work-in-progress push it
%% corresponds to.  A PDF with no commit hash cannot be checked against the
%% record and should be treated as a draft opinion, not a result."
%%
%% Written 2026-10-16 because the fault had recurred twice and I had twice
%% fixed only the instance:
%%   - 2026-10-14: the Q473 PDF went to Rick with no recipient and no hash.
%%   - 2026-10-14: the Q460 PDF went to Rick saying "Recipient: Robin Langer".
%%     He noticed, and took the commit hash as authority instead.
%% Both times I repaired the one file and wrote "fix the template, not just the
%% instance" in my notes.  This file is that fix.  \input it; do not retype it.
%%
%% USAGE, immediately after \maketitle:
%%   %% protocol-header.tex -- PROTOCOL §2.3 page-1 block, as a TEMPLATE not an instance.
%%
%% PROTOCOL §2.3: "Every PDF must state, on its first page: author, date,
%% recipient, a title, and the commit hash of the work-in-progress push it
%% corresponds to.  A PDF with no commit hash cannot be checked against the
%% record and should be treated as a draft opinion, not a result."
%%
%% Written 2026-10-16 because the fault had recurred twice and I had twice
%% fixed only the instance:
%%   - 2026-10-14: the Q473 PDF went to Rick with no recipient and no hash.
%%   - 2026-10-14: the Q460 PDF went to Rick saying "Recipient: Robin Langer".
%%     He noticed, and took the commit hash as authority instead.
%% Both times I repaired the one file and wrote "fix the template, not just the
%% instance" in my notes.  This file is that fix.  \input it; do not retype it.
%%
%% USAGE, immediately after \maketitle:
%%   %% protocol-header.tex -- PROTOCOL §2.3 page-1 block, as a TEMPLATE not an instance.
%%
%% PROTOCOL §2.3: "Every PDF must state, on its first page: author, date,
%% recipient, a title, and the commit hash of the work-in-progress push it
%% corresponds to.  A PDF with no commit hash cannot be checked against the
%% record and should be treated as a draft opinion, not a result."
%%
%% Written 2026-10-16 because the fault had recurred twice and I had twice
%% fixed only the instance:
%%   - 2026-10-14: the Q473 PDF went to Rick with no recipient and no hash.
%%   - 2026-10-14: the Q460 PDF went to Rick saying "Recipient: Robin Langer".
%%     He noticed, and took the commit hash as authority instead.
%% Both times I repaired the one file and wrote "fix the template, not just the
%% instance" in my notes.  This file is that fix.  \input it; do not retype it.
%%
%% USAGE, immediately after \maketitle:
%%   \input{protocol-header}
%%   \protocolheader{Rick (grandpa-rick)}{clio-vega/proofs}{60e4ce3}
%%
%% The recipient argument is MANDATORY and has no default on purpose: a default
%% is what produced "Recipient: Robin Langer" on a PDF addressed to Rick.

%% \maketitle CLEARS \@date (and \@author, \@title) in the article class, so a
%% macro called AFTER \maketitle that reads \@date prints an EMPTY row.  Caught
%% 2026-10-16 only because I ran pdftotext on page 1 instead of trusting rc=0:
%% the build was clean and the Date field was blank.  Hence: snapshot the date
%% HERE, at \input time in the preamble, where \@date is still live.
\let\protocol@date\@date

\newcommand{\protocolheader}[3]{%
  \begin{center}
  \fbox{\begin{minipage}{0.92\textwidth}
  \small
  \textbf{PROTOCOL §2.3 record.}\\[1mm]
  \begin{tabular}{@{}ll@{}}
  Author     & Clio \\
  Date       & \protocol@date \\
  Recipient  & #1 \\
  Repository & \texttt{#2} \\
  Commit     & \texttt{#3} \\
  \end{tabular}\\[1mm]
  This PDF corresponds to the pushed commit above.  Where this text and the
  commit disagree, \textbf{the commit is authoritative}.
  \end{minipage}}
  \end{center}
  \vspace{2mm}%
}

%%   \protocolheader{Rick (grandpa-rick)}{clio-vega/proofs}{60e4ce3}
%%
%% The recipient argument is MANDATORY and has no default on purpose: a default
%% is what produced "Recipient: Robin Langer" on a PDF addressed to Rick.

%% \maketitle CLEARS \@date (and \@author, \@title) in the article class, so a
%% macro called AFTER \maketitle that reads \@date prints an EMPTY row.  Caught
%% 2026-10-16 only because I ran pdftotext on page 1 instead of trusting rc=0:
%% the build was clean and the Date field was blank.  Hence: snapshot the date
%% HERE, at \input time in the preamble, where \@date is still live.
\let\protocol@date\@date

\newcommand{\protocolheader}[3]{%
  \begin{center}
  \fbox{\begin{minipage}{0.92\textwidth}
  \small
  \textbf{PROTOCOL §2.3 record.}\\[1mm]
  \begin{tabular}{@{}ll@{}}
  Author     & Clio \\
  Date       & \protocol@date \\
  Recipient  & #1 \\
  Repository & \texttt{#2} \\
  Commit     & \texttt{#3} \\
  \end{tabular}\\[1mm]
  This PDF corresponds to the pushed commit above.  Where this text and the
  commit disagree, \textbf{the commit is authoritative}.
  \end{minipage}}
  \end{center}
  \vspace{2mm}%
}

%%   \protocolheader{Rick (grandpa-rick)}{clio-vega/proofs}{60e4ce3}
%%
%% The recipient argument is MANDATORY and has no default on purpose: a default
%% is what produced "Recipient: Robin Langer" on a PDF addressed to Rick.

%% \maketitle CLEARS \@date (and \@author, \@title) in the article class, so a
%% macro called AFTER \maketitle that reads \@date prints an EMPTY row.  Caught
%% 2026-10-16 only because I ran pdftotext on page 1 instead of trusting rc=0:
%% the build was clean and the Date field was blank.  Hence: snapshot the date
%% HERE, at \input time in the preamble, where \@date is still live.
\let\protocol@date\@date

\newcommand{\protocolheader}[3]{%
  \begin{center}
  \fbox{\begin{minipage}{0.92\textwidth}
  \small
  \textbf{PROTOCOL §2.3 record.}\\[1mm]
  \begin{tabular}{@{}ll@{}}
  Author     & Clio \\
  Date       & \protocol@date \\
  Recipient  & #1 \\
  Repository & \texttt{#2} \\
  Commit     & \texttt{#3} \\
  \end{tabular}\\[1mm]
  This PDF corresponds to the pushed commit above.  Where this text and the
  commit disagree, \textbf{the commit is authoritative}.
  \end{minipage}}
  \end{center}
  \vspace{2mm}%
}

\makeatother

\newtheorem{thm}{Theorem}
\newtheorem{prop}[thm]{Proposition}
\newtheorem{lem}[thm]{Lemma}
\newtheorem{cor}[thm]{Corollary}
\theoremstyle{definition}
\newtheorem{defn}[thm]{Definition}
\theoremstyle{remark}
\newtheorem{rem}[thm]{Remark}

\DeclareMathOperator{\Inv}{Inv}
\DeclareMathOperator{\ord}{ord}
\newcommand{\At}{\widetilde{A}}
\newcommand{\bI}{\mathbf{I}}
\newcommand{\bM}{\mathbf{M}}
\newcommand{\sm}{\setminus}

\title{The boundary diagonal is a sum of two complementary counts\\
\normalsize $\ord_{S^{\varepsilon}_m}\At^{X}_{\bI,\bI}
 =\#\{a\in M_m:\;q^{[\,m\notin\bI\,]}\varepsilon^{a}=1\}$,\\
\normalsize and why a colength cannot tell \textup{(BD)} from its complement}
\author{Clio}
\date{2026-10-22}

\begin{document}
\maketitle
\protocolheader{Rick (grandpa-rick)}{clio-vega/proofs}{PENDING}

\begin{abstract}
I answer Q544 and Q546.  Both answers come from one theorem, which strictly
contains the master formula of \texttt{2026-10-21-c2} (\texttt{thm:BD}) and has
\emph{no hypothesis on $q$}: for every specialised $q\in\mathbb{C}^{\times}$ and every
$\varepsilon\in\mathbb{C}^{\times}$,
\[
  \ord_{S^{\varepsilon}_m}\At^{X}_{\bI,\bI}
  \;=\;[\,m\in\bI\,]\cdot f_{M}(\varepsilon)\;+\;[\,m\notin\bI\,]\cdot g_{M}(q,\varepsilon),
  \qquad
  \begin{aligned}
   f_{M}(\varepsilon)&:=\#\{a\in M: \varepsilon^{a}=1\},\\
   g_{M}(q,\varepsilon)&:=\#\{a\in M: q\,\varepsilon^{a}=1\},
  \end{aligned}
\]
where $M=M_m(X):=\{a>0: a\,t_m\in\Delta^{+}(X_N)\}$ is the \emph{multiplicity set} of the
boundary direction.  \texttt{lem:T}'s hypothesis $q\ne\varepsilon^{-h}$ \emph{is} the
statement that the second term vanishes; the composition point of Q544 is therefore not
the formula's breakdown but the point where its suppressed term is maximal.

\textbf{Q544: the two doublings do compose numerically, and the mechanism is a sum.}
At $(BC_N,\varepsilon=1,q=1)$ the order is $2$ for \emph{every} $\bI$ and the colength is
$2^{N+1}$, so the pre-registered multiplicative guess is numerically correct.  It is
mechanistically wrong: the colength is $2^{N-1}(f_M+g_M)$, and the guess $2\times2$ agrees
with the truth $2+2$ only by the arithmetic of $2$.  Neither ``doubling'' is a doubling.
$f$ is \emph{additive} in $M$, so $f_{BC}=f_B+f_C$ and not $2f_B$ --- at $\varepsilon=-1$
that reads $1=0+1\neq2\cdot0$; and the $q$-degeneration leaves the $m\in\bI$ half inert
while switching on a complementary term.  The first outright failure of multiplicativity
is $M=\{1,2,4\}$ ($3$, against a guess of $4$), covered by the theorem and realised by no
root system.

\textbf{A sharper consequence, pre-registered as F1 and confirmed.}  At
$(B_N,\varepsilon=-1,q=-1)$ the order is $[\,m\notin\bI\,]$ --- the exact
\emph{complement} of (BD) --- with colength $2^{N-1}$, the \emph{same integer} as (BD)
itself.  So \textbf{the colength cannot distinguish (BD) from its complement}: a
flattening and an inversion are the same number.

\textbf{Q546: false as asked, and the corrected form is a Möbius inversion.}  The
$\varepsilon$-profile $f_M$ at a primitive $d$-th root of unity is $\#(M\cap d\mathbb{Z})$,
a sum-over-multiples transform of $\mathbf 1_M$, so $[\,n\in M\,]=\sum_{k\ge1}\mu(k)f_M(kn)$:
the profile determines $M_m$ \emph{exactly and nothing more}.  $M_m$ is rank-blind, so
$B_2$ and $B_3$ have equal profiles and the literal falsifier is satisfied.  The
\emph{two}-point profile $(\varepsilon=\pm1)$ suffices for root systems only because
$M_m\subseteq\{1,2\}$ for every crystallographic system.
\end{abstract}

%=====================================================================
\section{What was open}

\texttt{2026-10-21-c2-G3-boundary-diagonal-typeB.tex} proves (\texttt{thm:BD}) that for
$X\in\{B,C\}$,
\begin{equation}
  \ord_{S^{\varepsilon}_m}\At^{X}_{\bI,\bI}
  \;=\;[\,m\in\bI\,]\cdot
  \#\{\,a>0:\;a\,t_m\in\Delta^{+}(X_N),\ \varepsilon^{a}=1\,\},
  \label{eq:old}
\end{equation}
under the hypothesis $q\ne\varepsilon^{-h}$ of its \texttt{lem:T}, which is automatic
because $q$ is transcendental over $\mathbb{C}$ in $F=\mathbb{C}(q)(w_1,\dots,w_N)$.  Two
questions were left open by that paper and sharpened by \textsc{dream} 10-21 c2.

\textbf{Q544.} Two separate ``doublings'' of the boundary colength were recorded.  A
non-reduced ambient ($BC_N$, in which $t_m$ \emph{and} $2t_m$ are positive roots) sends
the root count $1\to2$, measured $49/49$ at $N\le4$.  Independently, $q\to1$ at
$\varepsilon=1$ destroys the $\bI$-dependence of \eqref{eq:old} --- the $m\in\bI$ half is
inert ($1\to1$) and the $m\notin\bI$ half breaks ($0\to1$), measured $129/129$, both types,
$N\le5$ --- taking the colength $2^{N-1}\to2^{N}$.  Do they compose?  The point
$(BC_N,\varepsilon=1,q=1)$ is $q=1=\varepsilon^{-h}$, \emph{exactly} where \texttt{lem:T}
fails, so \eqref{eq:old} cannot be evaluated there and the natural guess $2^{N+1}$ had no
derivation behind it.

\textbf{Q546.} The measured pairs $(\ord$ at $\varepsilon=1,\ \ord$ at $\varepsilon=-1)$
are $(1,0)$, $(1,1)$, $(2,1)$ for $B$, $C$, $BC$ and separate all three, while no single
$\varepsilon$ does.  Is the pair a complete invariant of the ambient root system?

Both are answered by removing the hypothesis instead of working around it.

%=====================================================================
\section{Setup: $q$ becomes a scalar}

I keep every convention of the 10-21 paper except one: $q$ is specialised.

Write $M_m(X):=\{a>0: a\,t_m\in\Delta^{+}(X_N)\}$, the \emph{multiplicity set} of the
boundary direction $t_m$.  For the reduced systems $M_m(B)=\{1\}$ and $M_m(C)=\{2\}$; for
the non-reduced $BC_N$, $M_m(BC)=\{1,2\}$.  $M_m$ does not depend on $m$, and I drop the
subscript when convenient.

\begin{defn}[the specialised valuation]\label{def:ord}
Fix $q_0\in\mathbb{C}^{\times}$ and set $F:=\mathbb{C}(w_1,\dots,w_N)$,
$R:=\mathbb{C}[w^{\pm}]$.  Fix $m\in[N]$ and $\varepsilon\in\mathbb{C}^{\times}$.  Since
$R/(w_m-\varepsilon)\cong\mathbb{C}[w_j^{\pm}:j\ne m]$ is a domain, $(w_m-\varepsilon)$ is
a prime of $R$; let $\mathcal{O}\subset F$ be the local ring, $t:=w_m-\varepsilon$ the
uniformiser, $k=\mathcal{O}/(t)\cong\mathbb{C}(w_j:j\ne m)$ the residue field, and
$\ord_S$ the associated divisorial valuation on $F^{\times}$.
\end{defn}

The only change from \texttt{def:ord} of 10-21 is that the coefficient field is
$\mathbb{C}$ rather than $\mathbb{C}(q)$.  Nothing in the construction used the
transcendence of $q$: it used only that the $w_j$ with $j\ne m$ are algebraically
independent over the coefficient field, which they are over $\mathbb{C}$.

\begin{lem}[unit criterion, over $\mathbb{C}$]\label{lem:unit}
Let $M=\prod_j w_j^{a_j}$ be a Laurent monomial and $c\in\mathbb{C}^{\times}$.  Then
$1-cM\in\mathcal{O}$, and
\begin{enumerate}[label=\textup{(\roman*)},leftmargin=2.2em]
\item $\ord_S(1-cM)\ge1 \iff a_j=0$ for all $j\ne m$, and $c\,\varepsilon^{a_m}=1$;
\item if moreover $a:=a_m\ne0$ then $\ord_S(1-cM)=1$ exactly, with cofactor
  $a\,\varepsilon^{a-1}$ at the generic point of $S$.
\end{enumerate}
\end{lem}

\begin{proof}
Verbatim the proof of \texttt{lem:unit} of 10-21 with $\mathbb{C}(q)$ replaced by
$\mathbb{C}$.  Each $w_j$ is a unit of $\mathcal{O}$ ($\varepsilon\ne0$ for $j=m$), so
$1-cM\in\mathcal{O}$, with residue $1-c\,\varepsilon^{a_m}\prod_{j\ne m}w_j^{a_j}$ in $k$.
The $w_j$, $j\ne m$, are algebraically independent over $\mathbb{C}$, so this vanishes iff
the monomial is constant and the scalar is $1$, which is (i).  For (ii), (i) gives
$1-cM=1-\varepsilon^{-a}w_m^{a}$; for $a>0$,
$1-\varepsilon^{-a}w_m^{a}=-\varepsilon^{-a}(w_m-\varepsilon)\sum_{i=0}^{a-1}w_m^{i}\varepsilon^{a-1-i}$,
and the sum has residue $a\varepsilon^{a-1}\ne0$ in $k$ because
$\operatorname{char}\mathbb{C}=0$ and $\varepsilon\ne0$; for $a<0$ factor out the unit
$-\varepsilon^{-a}w_m^{a}$ and apply the case $-a>0$.
\end{proof}

The single substantive point is that \textbf{$c$ now ranges over a field containing
$q_0$}.  In the 10-21 argument $c\in\{1\}\cup q\mathbb{C}^{\times}$ with $q$
transcendental, so $c\varepsilon^{a}=1$ was \emph{unsatisfiable} whenever $c$ carried a
$q$; that unsatisfiability is the whole content of \texttt{lem:T}, and it is exactly what
specialisation removes.

%=====================================================================
\section{The theorem}

Recall from \texttt{eq:inv}, \texttt{eq:T}, \texttt{eq:diag} of the 10-21 paper that for
$\bI\subseteq[N]$ with strict partition $\lambda=\lambda(\bI)$,
\begin{align}
  \Inv(\lambda)&=\{a\,t_k : k\in\bI,\ a\in M\}
   \cup\{t_j+t_k: j<k,\ j,k\in\bI\}
   \cup\{t_k-t_j: j<k,\ j\notin\bI,\ k\in\bI\},
  \label{eq:inv}\\
  T(\lambda)&=\{a\,t_k : k\notin\bI,\ a\in M\}
   \cup\{t_j+t_k: j<k,\ j,k\notin\bI\}
   \cup\{t_k-t_j: j<k,\ j\in\bI,\ k\notin\bI\},
  \label{eq:T}\\
  \At^{X}_{\bI,\bI}
  &=\prod_{\alpha\in\Inv(\lambda)}\bigl(1-e^{-\alpha}\bigr)
   \prod_{\alpha\in T(\lambda)}\bigl(1-q\,e^{+\alpha}\bigr).
  \label{eq:diag}
\end{align}
For $X\in\{B,C\}$, $M$ is a singleton and \eqref{eq:inv}--\eqref{eq:diag} are Motegi's
theorem, read first-hand at L1305--1335 (type $B$) and L998--1046 (type $C$) of
\texttt{arXiv:2610.15344}.  For $X=BC$ there is no weight function in the literature:
\eqref{eq:diag} with $M=\{1,2\}$ is a \emph{definition}, the non-reduced extension
introduced as a lesion in \texttt{cor:red} of 10-21.  Theorem~\ref{thm:main} is therefore
a theorem about Motegi's object when $X\in\{B,C\}$ and about my extension when $X=BC$; I
flag this because the Q544 answer lives in the extension and only its mechanism is
corroborated in the literature's types.

\begin{thm}[boundary diagonal at specialised $q$]\label{thm:main}
Let $N\ge1$, $\bI\subseteq[N]$, $m\in[N]$, $M=M_m(X)$, and let
$q_0,\varepsilon\in\mathbb{C}^{\times}$ be arbitrary.  Then, with $q:=q_0$ in
\eqref{eq:diag},
\begin{equation}
  \ord_{S^{\varepsilon}_m}\At^{X}_{\bI,\bI}
  \;=\;\#\bigl\{\,a\in M \;:\; q_0^{\,[\,m\notin\bI\,]}\,\varepsilon^{a}=1\,\bigr\}
  \;=\;[\,m\in\bI\,]\,f_M(\varepsilon)+[\,m\notin\bI\,]\,g_M(q_0,\varepsilon),
  \label{eq:main}
\end{equation}
where $f_M(\varepsilon)=\#\{a\in M:\varepsilon^{a}=1\}$ and
$g_M(q_0,\varepsilon)=\#\{a\in M: q_0\varepsilon^{a}=1\}$.  No hypothesis on $q_0$ is
used.
\end{thm}

\begin{proof}
$\ord_S$ is a valuation, so by \eqref{eq:diag} it is the sum of the orders of the finitely
many factors, each of which lies in $\mathcal{O}$ by Lemma~\ref{lem:unit}.  Sort the
factors by the support of the exponent vector of the monomial they contain.

\emph{Pair and difference factors are units, whatever $q_0$ is.}  For
$\alpha=t_j\pm t_k$ with $j\ne k$ the exponent vector of $e^{\pm\alpha}$ is supported on
two \emph{distinct} indices, so some index $\ne m$ carries a nonzero exponent and
Lemma~\ref{lem:unit}(i) fails.  This argument never inspects $c$, hence is untouched by
specialisation.  The same applies to boundary factors $a\,t_k$ with $k\ne m$, whose
support is $\{k\}\not\ni m$.

\emph{The boundary factors at index $m$.}  By \eqref{eq:inv}--\eqref{eq:T}, for each
$a\in M$ the root $a\,t_m$ lies in $\Inv(\lambda)$ if $m\in\bI$ and in $T(\lambda)$ if
$m\notin\bI$ --- in exactly one of the two, and the choice is the same for every $a\in M$.
So the factors at index $m$ are
\[
  \{\,1-w_m^{-a}\;:\;a\in M\,\}\quad(m\in\bI),
  \qquad
  \{\,1-q_0\,w_m^{a}\;:\;a\in M\,\}\quad(m\notin\bI),
\]
one factor per $a\in M$.  Lemma~\ref{lem:unit}(i) applies with support $\{m\}$: the first
has $c=1$, $a_m=-a$, so its order is $\ge1$ iff $\varepsilon^{-a}=1$, i.e.\ iff
$\varepsilon^{a}=1$; the second has $c=q_0$, $a_m=a$, so its order is $\ge1$ iff
$q_0\varepsilon^{a}=1$.  In either case $a_m=\pm a\ne0$, so Lemma~\ref{lem:unit}(ii) makes
the order \emph{exactly} $1$ whenever it is positive.  Summing over $a\in M$ gives
$f_M(\varepsilon)$ in the first case and $g_M(q_0,\varepsilon)$ in the second, which is
\eqref{eq:main}.
\end{proof}

\begin{rem}[\texttt{thm:BD} is the locus where the second term dies]
\label{rem:recovers}
If $q$ is transcendental then $g_M(q,\varepsilon)=0$, because $q\varepsilon^{a}=1$ would
make $q$ algebraic; and $f_M(\varepsilon)=\#\{a>0: a\,t_m\in\Delta^{+},\varepsilon^{a}=1\}$
by definition of $M$.  So \eqref{eq:main} specialises to \eqref{eq:old}.  More precisely,
\texttt{lem:T}'s hypothesis $q\ne\varepsilon^{-h}$ is, for singleton $M=\{h\}$,
\emph{precisely} the condition $g_M(q,\varepsilon)=0$.  The hypothesis was not protecting
the theorem from a degeneration; it was naming the locus on which one of its two terms
vanishes.  The 10-21 lesion (``$q:=1$ leaves the $m\in\bI$ half inert and breaks the
$m\notin\bI$ half, $129/129$'') is \eqref{eq:main} read at $q_0=1$: $f$ does not mention
$q_0$ at all, and $g_M(1,\varepsilon)=f_M(\varepsilon)$.
\end{rem}

\begin{rem}[the geometry: two families of hyperplanes, and the indicator chooses one]
For each $a\in M$ the two boundary factors at index $m$ have zero divisors
\[
  \{\,w_m^{a}=1\,\}=\textstyle\sum_{\zeta\in\mu_a}\{w_m=\zeta\},
  \qquad
  \{\,q_0w_m^{a}=1\,\}=\textstyle\sum_{\eta^{a}=q_0^{-1}}\{w_m=\eta\},
\]
the $a$-torsion of the $w_m$-line and a coset of it.  All zeros are simple
(Lemma~\ref{lem:unit}(ii)).  So there are $2\,|M|$ candidate loci through index $m$; the
indicator $[\,m\in\bI\,]$ selects a half of them, and the order counts how many of the
selected ones pass through $\varepsilon$.  The coset equals the torsion iff $q_0=1$, which
is why $q_0=1$ makes the two halves agree for \emph{every} $\varepsilon$ and every $X$.
\end{rem}

%=====================================================================
\section{Q544: the composition point, and what ``doubling'' meant}

\begin{cor}[the colength is a sum]\label{cor:colength}
For every $q_0,\varepsilon\in\mathbb{C}^{\times}$,
\[
  \ord_{S^{\varepsilon}_m}\det\Phi=\sum_{\bI\subseteq[N]}\ord_S\At^{X}_{\bI,\bI}
  \;=\;2^{N-1}\bigl(f_M(\varepsilon)+g_M(q_0,\varepsilon)\bigr).
\]
\end{cor}
\begin{proof}
Exactly $2^{N-1}$ subsets contain $m$ and $2^{N-1}$ do not; apply \eqref{eq:main}.
\end{proof}

\begin{cor}[Q544]\label{cor:q544}
At $X=BC_N$, $\varepsilon=1$, $q_0=1$: $M=\{1,2\}$ and $f_M(1)=g_M(1,1)=2$, so
$\ord_S\At^{BC}_{\bI,\bI}=2$ for \emph{every} $\bI$ --- the order is constant in $\bI$ ---
and $\ord_S\det\Phi=2^{N-1}\cdot4=2^{N+1}$.
\end{cor}

So the pre-registered multiplicative prediction is \emph{numerically correct}, and the
reason is worth stating precisely, because it is not multiplicativity.

\begin{cor}[neither ``doubling'' is a doubling]\label{cor:nodouble}
\begin{enumerate}[label=\textup{(\alph*)},leftmargin=2.2em]
\item $f_M$ is \textbf{additive in $M$}: $f_{M_1\sqcup M_2}=f_{M_1}+f_{M_2}$, immediately
  from the definition.  Since $M(BC)=M(B)\sqcup M(C)$,
  \[
    f_{BC}(\varepsilon)=f_B(\varepsilon)+f_C(\varepsilon)\qquad\text{and not}\qquad
    2f_B(\varepsilon).
  \]
  The two differ: at $\varepsilon=-1$, $f_B=0$, $f_C=1$, $f_{BC}=1=0+1\ne2\cdot0$.  The
  whole measured $\{B,C,BC\}\times\{1,-1\}$ table is this one line --- the $BC$ row is the
  $B$ row plus the $C$ row --- and ``non-reducedness doubles the count'' is true only at
  $\varepsilon=1$, where $f_B=f_C$ by accident.
\item The $q$-degeneration does not multiply the order by anything.  It leaves
  $f_M(\varepsilon)$, the $m\in\bI$ half, \emph{untouched} ($f$ does not mention $q_0$),
  and switches on the complementary term $g_M$.  At $q_0=1$ that term happens to equal
  $f_M$, which is why the \emph{colength} doubles --- $2^{N-1}(f+f)=2\cdot2^{N-1}f$ --- and
  why it looked like a property of the order.
\item Hence the composition law is $f+g$, and the agreement with the multiplicative guess
  at Corollary~\ref{cor:q544} is the coincidence $2\cdot2=2+2$.  The first outright failure
  is $M=\{1,2,4\}$ at $\varepsilon=-1$: the generic-$q$ value is $f=2$, the multiplicative
  guess for $q_0=-1$ is $4$, and the truth is $f+g=2+1=3$.  No root system realises
  $M=\{1,2,4\}$, so within the crystallographic family multiplicativity holds --- for an
  arithmetic reason about $|M|\le2$, not a structural one.
\end{enumerate}
\end{cor}

\begin{cor}[F1: the indicator inverts, and the colength cannot see it]\label{cor:inv}
At $X=B_N$, $\varepsilon=-1$, $q_0=-1$: $M=\{1\}$, $f_M(-1)=0$ and $g_M(-1,-1)=1$, so
\[
  \ord_{S^{-1}_m}\At^{B}_{\bI,\bI}=[\,m\notin\bI\,],
\]
the exact \emph{complement} of \textup{(BD)}, with colength $2^{N-1}$ --- the \emph{same
integer} as (BD) itself.  The same inversion occurs at $(C_N,\varepsilon=i,q_0=-1)$ and at
$(BC_N,\varepsilon=i,q_0=-i)$.
\end{cor}

This is the sharp form of the Q544 finding, and it is the one I did not expect.  The 10-21
paper's closing warning reads ``a $q\to1$ limit destroys (BD)'s dependence on $\bI$
entirely''.  That is true at $q_0=1$ and is not the general phenomenon.  The general
phenomenon is that a second divisor family exists, is invisible while $q$ is
transcendental, and meets $S^{\varepsilon}_m$ on a locus of its own; depending on where
$(q_0,\varepsilon)$ sits relative to the two families, the order is $[m\in\bI]$,
$[m\notin\bI]$, constant, or zero.  \emph{Flattening is one case and inversion is another,
and the colength $2^{N-1}$ is the same number for (BD) and for its complement.}  Measured
at $N=3$, $m=2$: both have colength $4$, with supports
$\{\bI\ni2\}=\{\{2\},\{1,2\},\{2,3\},\{1,2,3\}\}$ and
$\{\bI\not\ni2\}=\{\varnothing,\{1\},\{3\},\{1,3\}\}$ respectively.

\begin{rem}[what this says about reporting]
Every statement of the form ``the boundary colength of my lattice is $2^{N-1}$'' is a
projection that forgets two things, not one: the divisor parameter $\varepsilon$
(Q546's point) \emph{and} the support of the order as a function of $\bI$.  The second
loss is the worse one, because the support is what (BD) asserts.  A report should carry
the pair $(f_M,g_M)$, from which the order, the support and the colength all follow.
\end{rem}

%=====================================================================
\section{Q546: the profile determines the multiplicity set, and nothing more}

\begin{thm}[the $\varepsilon$-profile is a Möbius transform of $\mathbf 1_M$]
\label{thm:profile}
For $d\ge1$ let $\varepsilon_d$ be any primitive $d$-th root of unity and
$f_M(d):=f_M(\varepsilon_d)$ --- well defined, since $\varepsilon_d^{a}=1\iff d\mid a$
depends only on $d$.  Then
\[
  f_M(d)=\#\bigl(M\cap d\mathbb{Z}\bigr)=\sum_{n\,:\,d\mid n}\mathbf 1_M(n),
  \qquad\text{hence}\qquad
  \mathbf 1_M(n)=\sum_{k\ge1}\mu(k)\,f_M(kn)\quad(n\ge1),
\]
the sum being finite because $M$ is finite.  Consequently the profile
$d\mapsto f_M(d)$ and the set $M$ determine one another.  For $\varepsilon$ not a root of
unity, $f_M(\varepsilon)=0$, so the profile carries no further information.
\end{thm}
\begin{proof}
The first display is the definition of $f_M$ together with $\varepsilon_d^{a}=1\iff d\mid a$.
The second is Möbius inversion for the sum-over-multiples transform: for $F(d)=\sum_{d\mid n}G(n)$
with $G$ finitely supported, $\sum_{k\ge1}\mu(k)F(kn)=\sum_{k\ge1}\mu(k)\sum_{kn\mid n'}G(n')
=\sum_{n\mid n'}G(n')\sum_{k\,\mid\,n'/n}\mu(k)=G(n)$, using
$\sum_{k\mid j}\mu(k)=[\,j=1\,]$.  Finally if $\varepsilon$ has infinite order then
$\varepsilon^{a}\ne1$ for all $a>0$, so $f_M(\varepsilon)=0$.
\end{proof}

\begin{cor}[Q546, as asked: \textbf{false}]\label{cor:q546no}
The $\varepsilon$-profile --- a fortiori the pair at $\varepsilon=\pm1$ --- is \emph{not} a
complete invariant of the ambient root system.  By Theorem~\ref{thm:profile} it is a
complete invariant of $M_m$, and $M_m$ is rank-blind: $M_m(B_2)=M_m(B_3)=\{1\}$, so $B_2$
and $B_3$ have identical profiles.  This is the pre-specified falsifier, and it is
satisfied.  The profile also cannot see the pair roots $t_k\pm t_j$ at all --- by the proof
of Theorem~\ref{thm:main} they are units for every $(q_0,\varepsilon)$ --- so no amount of
$\varepsilon$-data recovers them from this diagonal.
\end{cor}

\begin{cor}[Q546, corrected: \textbf{true}]\label{cor:q546yes}
The profile is a complete invariant of the multiplicity set $M_m(X)=\{a>0:a\,t_m\in\Delta^{+}\}$,
i.e.\ of the set of positive roots proportional to $t_m$.  Within
$\{B_N,C_N,BC_N\}$ at fixed $N$ the \emph{two}-point profile $(f_M(1),f_M(2))=(1,0),(1,1),(2,1)$
already separates the three, and this is not luck: $f_M(1)=|M|$ and $f_M(2)=\#\{a\in M: 2\mid a\}$
determine $M$ whenever $M\subseteq\{1,2\}$, and \textbf{every crystallographic root system has
$M_m\subseteq\{1,2\}$} --- the only non-reduced ones are of type $BC$, where $\alpha$ and
$2\alpha$ are roots and $3\alpha$ is not.  So the sufficiency of two points is a fact about
root systems, not about $\At$.  For $M\not\subseteq\{1,2\}$ two points no longer suffice:
$M=\{1\}$ and $M=\{3\}$ share $(f(1),f(2))=(1,0)$ and are separated first at $d=3$.
\end{cor}

\begin{rem}[what the full profile costs]
Reporting $(f_M(1),f_M(2))$ rather than one integer is the right repair \emph{for this
family}, but Corollary~\ref{cor:inv} shows it is still not enough on its own: the pair
$(f,g)$ --- not the profile of $f$ alone --- is what determines the order and its support.
At $(B_N,\varepsilon=-1,q_0=-1)$ the $f$-profile is $(f(1),f(2))=(1,0)$, identical to the
generic-$q$ type-$B$ reading, while the order is the complement of (BD).  The $f$-profile
answers ``which boundary roots resonate with $\varepsilon$''; it is silent about $q_0$.
\end{rem}

%=====================================================================
\section{Pre-registration: which arm fired}
\label{sec:prereg}

Written before computing, in \texttt{scratch/20261022-prove/PREREG.md}, and graded arm by
arm.  Seven falsifiers F1--F7 were pre-specified, \emph{six of them away from $q_0=1$},
because the pre-registration already noted that $q_0=1$ cannot discriminate:
there $g_M(1,\varepsilon)=f_M(\varepsilon)$ identically, so multiplicativity is vacuously
true at the very point Q544 names.

\begin{description}[leftmargin=1.4em]
\item[Q544 numbers.] Predicted $2$ on $m\in\bI$, $2$ on $m\notin\bI$, colength $2^{N+1}$,
  order constant in $\bI$: \textbf{all four TRUE} (Corollary~\ref{cor:q544}).
\item[The mechanism.] Predicted the two-term sum \eqref{eq:main} with
  $g=\#\{a\in M:q\varepsilon^{a}=1\}$ and ``$q\ne\varepsilon^{-h}$ is exactly $g=0$'':
  \textbf{TRUE} (Theorem~\ref{thm:main}, Remark~\ref{rem:recovers}).  This is the arm that
  matters: the numerical prediction and the mechanism prediction were registered
  separately, and only the second distinguishes a sum from a product.
\item[F1 (inversion in type $B$), F2 (type $C$), F3 ($BC$).] Predicted
  $\ord=[\,m\notin\bI\,]$: \textbf{all three TRUE}, $34/34$ each at $N\le3$ on the factor
  routes and, for F1 and F2, $34/34$ on the independent symmetrisation route.
\item[F4 ($BC$, $\varepsilon=1$, $q_0=-1$): $2[\,m\in\bI\,]$.] \textbf{TRUE}, $34/34$.
\item[F5 ($BC$, $\varepsilon=-1$, $q_0=1$): $\equiv1$.] \textbf{TRUE}, $34/34$.
\item[F6 (control, $B$, $\varepsilon=i$, $q_0=2$): $\equiv0$.] \textbf{TRUE}, $34/34$ ---
  a genuine control, since here the prediction is $0$ on both halves, so the formula is not
  fitted to the resonant points.
\item[F7 (the Q544 point itself): $\equiv2$.] \textbf{TRUE}, $34/34$.
\item[Q546.] Predicted ``\textbf{NO} as stated, \textbf{YES} in the corrected form'', with
  the literal falsifier satisfied by $B_2$ vs $B_3$ and the corrected statement a Möbius
  inversion over multiples: \textbf{TRUE on both halves}
  (Corollaries~\ref{cor:q546no}, \ref{cor:q546yes}).
\item[Not pre-registered, and the session's actual find.] That the colength is the
  \emph{same integer} $2^{N-1}$ for (BD) and for its complement
  (Corollary~\ref{cor:inv}).  I registered the inversion; I did not register that it is
  invisible to the quantity this programme reports.  The additivity $f_{BC}=f_B+f_C$
  (Corollary~\ref{cor:nodouble}(a)) was likewise not registered --- I had written
  ``non-reducedness doubles'' in the pre-registration's own framing of the prior state,
  and the $\varepsilon=-1$ column of my own measured table had been contradicting it since
  10-21.
\end{description}

%=====================================================================
\section{Verification}
\label{sec:verif}

Pure Python $+$ SymPy (\textbf{SageMath is not installed}).  $q_0$ and $\varepsilon$ are
exact SymPy numbers; $w_{j\ne m}$ stay \emph{symbolic}, so every valuation is taken at the
generic point of $S$ and the codimension-$2$ hypersurfaces discussed in
\texttt{rem:retract} of 10-21 cannot be hit.  Scripts:
\texttt{scratch/20261022-prove/\{q544.py, routeC.py, q546.py, colength.py\}}.

Three routes, of decreasing dependence on the objects under study.

\begin{description}[leftmargin=1.4em]
\item[Route A --- factorwise.]  $\ord$ computed on each factor of
  \eqref{eq:diag} after $q:=q_0$, and summed.  Uses additivity of the valuation and the
  factor list.
\item[Route B --- whole product.]  The factors are multiplied first, $q:=q_0$ substituted,
  and \emph{one} valuation taken.  Does not use additivity, so it is a check on the
  bookkeeping that Lemma~\ref{lem:unit} partitions.
\item[Route C --- symmetrisation.]  $\At^{X}_{\bI,\bI}$ computed from Motegi's
  $S_n$-symmetrisation \emph{definition} (\texttt{eq Motherfunction} L473--490, with
  $\mathcal{E}^{B}$ L1361--1367 / $\mathcal{E}^{C}$ L1057--1062), $q$ kept symbolic,
  \texttt{cancel}led, and \emph{then} specialised.  Uses neither \eqref{eq:diag} nor the
  root system --- the engine never mentions a positive root --- so it is independent of
  everything Theorem~\ref{thm:main} reasons about.  Available for $X\in\{B,C\}$ only; $BC$
  has no symmetrisation, being defined by \eqref{eq:diag}.
\end{description}

\paragraph{Domain check, pre-registered.}
The 10-19 lesson (``$0/18$ meant \emph{undefined}, not false'') was pre-registered here as
an explicit check, since $q_0=1$ is a degeneration.  Route C reduces
$\At^{X}_{\bI,\bI}$ to lowest terms as a rational function of $q$ and then divides out
every common factor of $(q-q_0)$ before substituting, reporting the count.  \textbf{The
count was $0$ in every one of the $7\times34$ cases, and neither numerator nor denominator
vanished at $q_0$ in any case.}  So $q_0=1$ is in the domain: $\At_{\bI,\bI}$ is regular
and nonzero there, and no $0/0$ is being read as a number.  This is as the theory predicts
--- \eqref{eq:diag} exhibits $\At_{\bI,\bI}$ as a \emph{polynomial} in $q$, hence regular
at every $q_0$ --- but the prediction and the measurement are independent.

\paragraph{Measurements.}
\begin{center}\small
\begin{tabular}{llll}
 arm & $(X,\varepsilon,q_0)$ & route A, $N\le5$ & route B, $N\le3$ & route C, $N\le3$\\\hline
 F7 \ Q544 point        & $(BC,1,1)$     & $258/258$ & $34/34$ & --- ($BC$)\\
 F1 \ inversion, type $B$ & $(B,-1,-1)$  & $258/258$ & $34/34$ & $34/34$\\
 F2 \ inversion, type $C$ & $(C,i,-1)$   & $258/258$ & $34/34$ & $34/34$\\
 F3 \ inversion, $BC$     & $(BC,i,-i)$  & $258/258$ & $34/34$ & --- ($BC$)\\
 F4 \ asymmetric counts   & $(BC,1,-1)$  & $258/258$ & $34/34$ & --- ($BC$)\\
 F5 \ $q_0=1$ off $\varepsilon=1$ & $(BC,-1,1)$ & $258/258$ & $34/34$ & --- ($BC$)\\
 F6 \ control, $\equiv0$  & $(B,i,2)$    & $258/258$ & $34/34$ & $34/34$\\
 \ \ \ $q_0=1$, type $B$  & $(B,1,1)$    & $258/258$ & $34/34$ & $34/34$\\
 \ \ \ $q_0=1$, type $C$  & $(C,-1,1)$   & $258/258$ & $34/34$ & $34/34$\\
 \ \ \ 10-21 baseline $B$ & $(B,1,3/7)$  & $258/258$ & $34/34$ & $34/34$\\
 \ \ \ 10-21 baseline $C$ & $(C,-1,3/7)$ & $258/258$ & $34/34$ & $34/34$\\
 \ \ \ 10-21 baseline $BC$& $(BC,1,3/7)$ & $258/258$ & $34/34$ & --- ($BC$)\\
\end{tabular}
\end{center}
The counts are all $(\bI,m)$ pairs: $\sum_{N\le3}2^{N}N=34$, and
$34+64+160=258$ for $N\le5$.  Route A reaches $N=5$ because each factor it inspects is a
binomial; routes B and C cancel a product of $\Theta(N^{2})$ factors and stop at $N=3$.
In every arm
the measured order was \emph{constant on each of the two halves} $m\in\bI$, $m\notin\bI$,
as \eqref{eq:main} requires; the two constants are $(f_M,g_M)$ and are printed in the
script output.  The three baseline rows reproduce the 10-21 numbers on the same domain, so
the specialised instrument agrees with the transcendental one where both apply.

\begin{description}[leftmargin=1.4em]
\item[Generalised multiplicity sets.]  \eqref{eq:main} measured for $M$ ranging over
  $\{1\},\{2\},\{1,2\},\{3\},\{1,3\},\{2,3\},\{1,2,3\},\{2,4\},\{1,2,3,4\},\{6\}$ ---
  six of which no root system realises --- against eight points
  $(\varepsilon,q_0)$ including $(\zeta_3,3/7)$, $(i,-1)$, $(i,-i)$,
  $(\zeta_3,\zeta_3^{-1})$: \textbf{$2720/2720$} at $N\le3$.  This is the measurement that
  bears on Theorem~\ref{thm:profile}, whose content is the map $M\mapsto f_M$.
\item[Colength.]  Corollary~\ref{cor:colength} measured by summing the order over all
  $\bI$: \textbf{$60/60$} over ten $(M,\varepsilon,q_0)$ points at $N\le3$, including
  $M=\{1,2,4\}$ at $(\!-1,-1)$, where the measured colength is $3\cdot2^{N-1}$ against the
  multiplicative guess $4\cdot2^{N-1}$ --- Corollary~\ref{cor:nodouble}(c).
\item[The support.]  At $N=3$, $m=2$: $(B,1,3/7)$ gives colength $4$ with support
  $\{\bI\ni2\}$; $(B,-1,-1)$ gives colength $4$ with support $\{\bI\not\ni2\}$.  Equal
  integer, complementary support.
\item[Möbius recovery.]  $\mathbf 1_M(n)=\sum_k\mu(k)f_M(kn)$ recovered $M$ exactly in
  $11/11$ cases (the ten above plus $\{4,6\}$), with the profile computed from
  \emph{explicit} primitive roots of unity $e^{2\pi i/d}$, $d\le12$, not from the
  divisibility rule --- so the identification $\varepsilon_d^{a}=1\iff d\mid a$ is
  measured rather than assumed.
\end{description}

\begin{rem}[instrument plumbing, recorded]
The trajectory emitter was initialised with an \emph{absolute} \texttt{log\_dir} and a
\emph{stable} session tag \texttt{Q544-20261022}, which is the 10-18 repair: the default
timestamp session fragmented one session into four files while \texttt{report} reads one.
A first call made before the repair wrote a bare-timestamp file; it was deleted and the
event re-emitted into the stable sink rather than left as a second fragment.
$18$ events emitted, $18$ present in the report --- counted, not assumed, since a
\texttt{verify} lacking \texttt{verification\_target} is silently dropped.

I also corrected one of my own encodings, unprompted by the checker.  The Q546 claim
listed the generalised-$M$ computation in its \texttt{inputs}, while the \texttt{verify}
event for that claim offered the \emph{same} computation as evidence --- so the
verification's evidence was an ancestor of its own target.  This is the 2026-10-18 fault
class (reading \texttt{inputs} as ``what this is \emph{about}'' rather than ``what this
\emph{depends on}''), and it is its fifth instance.  Theorem~\ref{thm:profile} was in fact
proved by hand from Theorem~\ref{thm:main}; the computation came afterwards and is
evidence, not a dependency.  Re-emitted with \texttt{inputs} corrected; the as-recorded
events are retained in the same file rather than deleted.
\end{rem}

%=====================================================================
\section{What this closes, and what it does not}

\begin{description}[leftmargin=1.4em]
\item[Q544 --- CLOSED.]  Corollary~\ref{cor:q544}: order $\equiv2$, colength $2^{N+1}$,
  constant in $\bI$.  The two doublings compose as a \emph{sum} $f+g$, and the
  multiplicative guess is right at this point only because $2\cdot2=2+2$
  (Corollary~\ref{cor:nodouble}).
\item[Q546 --- CLOSED, both halves.]  False as asked
  (Corollary~\ref{cor:q546no}); true as the statement that the profile determines $M_m$
  (Corollaries~\ref{cor:q546yes}, Theorem~\ref{thm:profile}).
\item[\texttt{thm:BD} --- strictly generalised, hypothesis removed.]
  Theorem~\ref{thm:main} has no hypothesis on $q$ and contains \eqref{eq:old} as the locus
  $g_M=0$.  Nothing in the 10-21 paper is retracted; one of its closing warnings is
  \emph{narrowed} --- ``a $q\to1$ limit destroys (BD)'s dependence on $\bI$'' is the
  $q_0=1$ case of a family that also inverts the dependence
  (Corollary~\ref{cor:inv}).
\item[A caveat on scope.]  For $X=BC$ the object is defined by \eqref{eq:diag}; route C
  does not apply and the $BC$ arms rest on routes A and B.  The \emph{mechanism} ---
  the two complementary blocks, and the inversion --- is confirmed on the independent
  symmetrisation route in types $B$ and $C$, which are Motegi's.
\item[Unchanged.]  \textbf{G2} (non-cancellation at fixed $\chi$) is untouched; nothing
  above bears on it.  The shared unverified RTV stable-envelope root beneath
  \eqref{eq:diag} is not reduced by this work --- Theorem~\ref{thm:main} takes
  \eqref{eq:diag} as its input exactly as 10-21 did.
\item[Downstream obligation.]  Every colength in \S\S5--7 of
  \texttt{2026-10-19-c2-sa-divisors-full-q.tex} is stated at transcendental $q$ and so sits
  at $g_M=0$; those statements stand.  What must not be inherited is the \emph{form}
  $2^{N-1}$: under any specialisation of $q$ the correct report is the pair
  $(f_M,g_M)$, which determines the order, its support, and the colength, and which no
  single integer determines.
\end{description}

\end{document}
